
\subsection{\texorpdfstring{\textbf{
      Обучение классической и непрерывной моделей Хопфилда}}
  {Обучение классической и непрерывной моделей Хопфилда}}

\textbf{Правило Хэбба для дискретной динамики.}

Для компонент паттернов \(p_{i}\) взятых из \(\ \{ 0,1\}\):


\[
  T_{ij} = \sum_{k}^{}{(2p_{i}^{k} - 1)(2p_{j}^{k} - 1)}
\]


Для \(\{ - 1,1\}\):


\[
  T_{ij} = \frac{1}{n}\sum_{k = 1}^{n}{p_{i}^{k}p_{j}^{k}}
\]


Важно, что сама матрица весов T непрерывна и
на диагонали всегда 0 - \(T_{ii}\) = 0.

При этом, Хэбб и Сторки обучают модель за одну итерацию.

\textbf{Сторки для непрерывной и дискретной динамики.}

\cite{storkey1997increasing}


\textbf{Скоростной градиент для непрерывной динамики}
\[
  \tau_f \frac{dx_i}{dt} =
  \sum\limits_{j=1}^{N_f} T_{ij} V_j - x_i + I_i
\]
рассматривается в виде динамической системы
\(\dot{x} = \omega(x,\theta,t)\), где \(\theta\) - веса матрицы
\(T\).

Скоростной градиент имеет вид
\[
  \frac{d\theta}{dt} = -\Gamma \nabla_{\theta}{\omega(x,\theta,t)}
\]
\[
  \omega(x,\theta,t) = \frac{dQ}{dt}
\]

Функционал \(Q\)

локальный \(Q = Q(x(t), t)\)

интегральный \(Q = \int\limits_0^t R(x(s),\theta(s),s) ds\) \cite{fradkov1990adaptive}

Можно заметить, что в нашем случае система адаптивная,
где устойчивость достигается за счёт настройки только
внутренних параметров самой системы.

% \cite{fradkov2003cybernetical_ru}